% TEMPLATE-MILC-v3.tex - plantilla maestra UNICA de las guias MILC v3.
%
% La usan las cuatro familias de guias, cada una con su driver:
%   · Grados 6-11          -> scripts/build-guias-g11.py
%   · Bebras               -> scripts/build-guias-bebras.py
%   · Semillero            -> scripts/build-guias-semillero.py
%   · Territorio interior  -> scripts/build-guias-territorio-interior.py
%
% Antes habia tres copias casi identicas (~400 lineas iguales, 6 distintas):
% cada mejora habia que aplicarla tres veces, y en la practica no se hacia
% (el motor de recursos vivia solo en dos de ellas). Lo que varia entre
% familias esta parametrizado en 6 placeholders que cada driver llena
% (se nombran aqui SIN su sintaxis real: el builder reemplaza por texto plano,
% tambien dentro de los comentarios, y un valor multilinea rompe el preambulo):
%
%   HEADER_IZQ        encabezado izquierdo de pagina
%   HEADER_DER        encabezado derecho de pagina
%   PORTADA_ETIQUETA  rotulo sobre el numero grande (GRADO/MOMENTO/MODULO)
%   PORTADA_SERIE     linea de serie de la portada
%   PORTADA_ID        identificador de la guia en la portada
%   PORTADA_CIRCULO   el circulito de grado (°), o vacio si no aplica
%   PDF_TITULO        titulo en texto plano (sin \textbf ni comillas TeX) para
%                     los metadatos del PDF (/Title); lo produce lib_xelatex.texto_pdf
%
% Composicion (auditoria 2026-09-03): las bandas de estacion piden espacio
% (\Needspace*) y prohiben el salto justo despues (\nopagebreak), asi no quedan
% huerfanas al pie; las cajas largas (softbox, drawbox, infoband, puentebox) son
% `breakable`, asi una caja de media pagina ya no arrastra todo a la siguiente
% dejando la anterior al 40 %. Todo texto sobre mostaza o verde va en milcNegro
% (blanco sobre #E5B400 da 1,93:1 y sobre #58A923 2,95:1; ninguno pasa WCAG AA).
% El resto de claves las documenta content/guias/_SCHEMA.md. El script valida
% que no queden placeholders sin reemplazar antes de escribir el archivo final.

% Etiquetado PDF (latex-lab): /Lang, estructura y texto alternativo de las
% figuras, para lectores de pantalla. Debe ir ANTES de \documentclass.
\DocumentMetadata{lang=es-CO,tagging=on}
\documentclass[11pt,letterpaper]{article}
% headheight=24pt: la cabecera derecha lleva dos lineas (identidad de la guia
% y estacion actual). headsep se acorta en la misma medida para que el bloque
% de texto quede exactamente donde estaba con headheight=12pt.
\usepackage[margin=2.0cm,headheight=24pt,headsep=13pt]{geometry}
\usepackage[table]{xcolor}
\usepackage{fontspec}
\usepackage[spanish]{babel}
\usepackage{tikz}
% Librerías TikZ para los diagramas de las guías (bloque `recursos:`):
%  · babel  → babel en español vuelve activos caracteres como «>», y TikZ los usa
%             en sus opciones (>=stealth, ->). Sin esta librería, cualquier
%             diagrama con flechas revienta con «\language@active@arg> has an
%             extra }». Debe ir ANTES de que se dibuje nada.
%  · positioning        → sintaxis «below left=2pt and 3pt».
%  · decorations.pathreplacing → llaves ({brace}) para señalar rangos.
\usetikzlibrary{babel,positioning,decorations.pathreplacing}
\usepackage{graphicx}
% Circuitos: el trabajo de electrónica se hace hoy con micro:bit y sus sensores
% integrados (MakeCode, sin cableado), así que no se necesitan esquemáticos.
% Si algún día entran montajes con protoboard, basta descomentar esta línea:
% los diagramas .tex/.tikz declarados en `recursos:` ya se incrustan solos.
% \usepackage{circuitikz}
\usepackage[most]{tcolorbox}
\usepackage{tabularx}
\usepackage{array}
\usepackage{fancyhdr}
\usepackage{titlesec}
\usepackage[document]{ragged2e}  % bandera a la derecha en todo el cuerpo
\usepackage{enumitem}
\usepackage{microtype}
\usepackage{etoolbox}
\usepackage{needspace}
% hyperref al final del preambulo: metadatos del PDF (titulo, autor, idioma,
% familia), marcadores por estacion (\pdfbookmark) y enlaces sin recuadro.
\usepackage[hidelinks,
  pdftitle={Visualización honesta — gráficos que mienten},
  pdfauthor={PhD. Álvaro Cárdenas Orozco},
  pdfsubject={Tecnología e Informática · Grado 9 · Guía 28 · MILC v3},
  pdflang={es-CO},
  bookmarksopen=true,bookmarksnumbered=false]{hyperref}

% Helvetica Neue no trae la flecha U+2192, asi que cada «→» escrito literalmente
% en el YAML se compilaba a NADA, en silencio (462 flechas en 61 guias: rutas del
% tipo «paleta → tipografia → lockup» salian rotas). Mapeamos el caracter a su
% version matematica, que si tiene glifo. Asi el autor puede seguir escribiendo
% «→» comodamente en el YAML.
\usepackage{newunicodechar}
\newunicodechar{→}{\ensuremath{\rightarrow}}
% Mismo problema con los simbolos de marca y vinetas: el «✓» de «una palomita ✓»
% salia como cuadrito vacio en el PDF de todas las guias que lo usaban.
\usepackage{pifont}
\newunicodechar{✓}{\ding{51}}
\newunicodechar{✗}{\ding{55}}
\newunicodechar{✔}{\ding{52}}
\newunicodechar{•}{\textbullet}
\newunicodechar{≥}{\ensuremath{\geq}}
\newunicodechar{≤}{\ensuremath{\leq}}

\setmainfont{Helvetica Neue}
\setsansfont{Helvetica Neue}
\setlength{\parindent}{0pt}
\setlength{\parskip}{6pt}
% Sin particion de palabras: en 6.º-7.º una palabra cortada por guion al final
% de linea («Comuni-cacion») frena la lectura mucho mas que una linea corta.
\hyphenpenalty=10000
\exhyphenpenalty=10000
\pretolerance=2000
\tolerance=3000
\setlength{\emergencystretch}{3em}
\linespread{1.10}
\renewcommand{\arraystretch}{1.25}
\setlist{leftmargin=5mm,itemsep=1mm,topsep=1mm}
\newcolumntype{Y}{>{\RaggedRight\arraybackslash}X}

\definecolor{milcMagenta}{HTML}{E6007E}
\definecolor{milcVino}{HTML}{5A0038}
\definecolor{milcTurquesa}{HTML}{0093A5}
\definecolor{milcVerde}{HTML}{58A923}
\definecolor{milcMostaza}{HTML}{E5B400}
\definecolor{milcOcre}{HTML}{B86600}
\definecolor{milcNegro}{HTML}{241816}
\definecolor{milcGris}{HTML}{F7F7F4}
\definecolor{milcLinea}{HTML}{E8E2DD}
\definecolor{milcGradePrimary}{HTML}{7C3AED}
\definecolor{milcGradeDark}{HTML}{4C1D95}
\definecolor{milcGradeSoft}{HTML}{EDE9FE}
\definecolor{milcGradeAccent}{HTML}{CFE7FF}
\definecolor{milcGradeText}{HTML}{FFFFFF}
\color{milcNegro}

\pagestyle{fancy}
\fancyhf{}
% Cabecera de dos lineas: identidad de la guia arriba y, a la derecha y en
% gris, la estacion en curso (\markright desde \milcestacion). Antes de la
% primera estacion la segunda linea va vacia; el \strut mantiene la altura
% para que la linea de arriba no baile entre paginas.
\lhead{\small Tecnología e Informática · Grado 9\\[-1pt]{\footnotesize\strut}}
\rhead{\small Guía 28 · MILC v3\\[-1pt]{\footnotesize\strut\color{milcNegro!65}\rightmark}}
\cfoot{\small\thepage}
\renewcommand{\headrulewidth}{0pt}

% Sobre mostaza (#E5B400) y verde (#58A923) el texto va en milcNegro; sobre el
% resto de la paleta, en blanco. Se usa dentro de fonttitle/fontupper, donde un
% \color posterior manda sobre coltitle.
\newcommand{\milctintasobre}[1]{%
  \ifboolexpr{test{\ifstrequal{#1}{milcMostaza}} or test{\ifstrequal{#1}{milcVerde}}}%
    {\color{milcNegro}}{\color{white}}}

\newtcolorbox{titlebox}[1]{
  enhanced,colback=#1,colframe=#1,arc=7mm,boxrule=0pt,
  left=7mm,right=7mm,top=3mm,bottom=3mm,
  before skip=8pt,after skip=8pt,
  fontupper=\sffamily\bfseries\Large\color{white}
}

\newtcolorbox{softbox}[3]{
  enhanced,breakable,pad at break=3mm,
  colback=#2,colframe=#1,borderline west={4pt}{0pt}{#1},
  arc=2mm,boxrule=.4pt,left=4mm,right=4mm,top=3mm,bottom=3mm,
  before skip=6pt,after skip=8pt,title={#3},coltitle=white,
  colbacktitle=#1,fonttitle=\sffamily\bfseries\milctintasobre{#1},fontupper=\RaggedRight,
  attach boxed title to top left={xshift=3mm,yshift=-2mm},
  boxed title style={arc=2mm, boxrule=0pt}
}

\newtcolorbox{drawbox}[2]{
  enhanced,breakable,pad at break=4mm,
  colback=white,colframe=#1,arc=4mm,boxrule=1pt,
  left=5mm,right=5mm,top=4mm,bottom=4mm,before skip=8pt,after skip=8pt,
  title={#2},coltitle=white,colbacktitle=#1,fonttitle=\sffamily\bfseries\milctintasobre{#1},
  fontupper=\RaggedRight,
  attach boxed title to top left={xshift=4mm,yshift=-2mm},
  boxed title style={arc=3mm, boxrule=0pt}
}

\newtcolorbox{infoband}[1]{
  enhanced,breakable,pad at break=2mm,colback=#1!8,colframe=#1!50,arc=2mm,boxrule=.5pt,
  left=4mm,right=4mm,top=2mm,bottom=2mm,before skip=4pt,after skip=4pt,
  fontupper=\small\RaggedRight
}

% Puente narrativo entre secciones (contrato editorial Conéctate).
% Da continuidad: "Acabas de X. Ahora vas a Y porque Z."
% Visualmente: banda izquierda en color del grado, texto en itálica pequeña.
\newtcolorbox{puentebox}{
  enhanced,breakable,pad at break=2mm,colback=milcGradeSoft,colframe=milcGradeDark,
  borderline west={3pt}{0pt}{milcGradeDark},
  arc=0mm,boxrule=0pt,left=5mm,right=4mm,top=2.5mm,bottom=2.5mm,
  before skip=4pt,after skip=8pt,
  fontupper=\itshape\small\color{milcNegro}\RaggedRight
}
% La flecha va en modo matematico a proposito: Helvetica Neue NO tiene el glifo
% U+2192, asi que \textrightarrow se compilaba a NADA («Missing character: There
% is no → in font Helvetica Neue») y el puente salia sin flecha. En modo
% matematico la flecha viene de la fuente matematica y si se dibuja.
\newcommand{\puente}[1]{%
  \begin{puentebox}%
    \textcolor{milcGradeDark}{$\rightarrow$}\hspace{2mm}#1%
  \end{puentebox}%
}

\newcommand{\linead}[1]{\noindent\color{milcLinea}\rule{#1}{0.5pt}\color{milcNegro}}
\newcommand{\respuesta}[1]{\par\vspace{2mm}\foreach \n in {1,...,#1}{\noindent\color{milcLinea}\rule{\linewidth}{0.5pt}\par\vspace{5mm}}\color{milcNegro}}
\newcommand{\checkbox}{\textcolor{milcGradeDark}{\fbox{\phantom{x}}}\hspace{5pt}}

% ─── Listas aireadas para las guías socioemocionales ────────────────────
% Componer las verificaciones como «\checkbox texto\\» deja el texto que salta
% de línea debajo del cuadrito y apelmaza el bloque. Estos entornos alinean la
% continuación y airean los ítems. Los usa Territorio Interior; cualquier
% familia puede hacerlo.
\newlist{checklista}{itemize}{1}
\setlist[checklista]{
  label=\textcolor{milcGradeDark}{\fbox{\phantom{x}}},
  leftmargin=9mm, labelsep=3.2mm, itemsep=2.4mm,
  topsep=2.5mm, parsep=0pt, partopsep=0pt, before=\RaggedRight
}
\newlist{pasoslista}{enumerate}{1}
\setlist[pasoslista]{
  label=\textcolor{milcGradeDark}{\sffamily\bfseries\arabic*.},
  leftmargin=8mm, labelsep=3mm, itemsep=2.8mm,
  topsep=1mm, parsep=0pt, partopsep=0pt, before=\RaggedRight
}

\newcommand{\phasecard}[4]{
\begin{tcolorbox}[enhanced,colback=#3,colframe=#2,arc=3mm,boxrule=.5pt,height=3.05cm,valign=center,halign=center,left=1.5mm,right=1.5mm,top=2mm,bottom=2mm]
{\sffamily\bfseries\fontsize{22}{24}\selectfont\textcolor{#2}{#1}}\par
{\sffamily\bfseries\small #4}
\end{tcolorbox}
}

% ── Piezas del rediseno 2026 (legibilidad 6.º-11.º) ───────────────────
% Banda de estacion: reemplaza el titlebox macizo. Titulo a la izquierda,
% dato practico (tiempo/modalidad) a la derecha, en una sola linea.
% Nunca huerfana: pide 9 lineas libres (o salta de pagina rellenando con
% \vfill) y prohibe el salto justo despues de la banda. Nueve y no seis:
% la caja partible que sigue necesita ~80pt (titulo adosado, relleno y dos
% lineas) para arrancar en la misma pagina; con menos, tcolorbox la manda
% entera a la siguiente y la banda se queda sola. El penalty 10000 va
% en `after`, ANTES del espacio posterior: un `after skip` (aunque sea 0pt)
% es un \vskip tras la caja, y ese glue es un punto de corte legal que un
% \nopagebreak escrito despues de \end{tcolorbox} ya no cierra. Cada banda
% deja marcador PDF y actualiza la cabecera derecha.
\newcounter{milcestacion}
\newcommand{\milcestacion}[2]{%
\Needspace*{9\baselineskip}%
\stepcounter{milcestacion}%
\pdfbookmark[1]{#2}{milcestacion\themilcestacion}%
\markright{#2}%
\begin{tcolorbox}[enhanced,colback=#1,colframe=#1,arc=2mm,boxrule=0pt,
  left=5mm,right=5mm,top=2.6mm,bottom=2.6mm,before skip=11pt,
  after={\par\nopagebreak\vskip7pt}]
{\sffamily\bfseries\fontsize{15}{18}\selectfont\milctintasobre{#1}#2}
\end{tcolorbox}}

% Tarjeta de paso numerada: sustituye las tablas «Pilar / Aplicacion», que
% eran formato de consulta docente, no de estudiante. El numero va en un
% circulo sobre el borde para que la secuencia se lea de un vistazo.
\newcommand{\milcpaso}[3]{%
\Needspace{4\baselineskip}%
\begin{tcolorbox}[enhanced,colback=white,colframe=milcLinea,arc=1.5mm,boxrule=.9pt,
  left=14mm,right=4mm,top=3mm,bottom=3mm,before skip=4pt,after skip=4pt,
  fontupper=\RaggedRight,
  overlay={\node[anchor=north west,circle,fill=#2,text=white,inner sep=0pt,
    minimum size=7.4mm,font=\sffamily\bfseries\small]
    at ([xshift=3.6mm,yshift=-3.2mm]frame.north west) {#1};}]
#3
\end{tcolorbox}}

% Casilla marcable de tamano util con lapiz (la anterior era \fbox{\phantom{x}},
% de alto variable segun el texto que la seguia).
\newcommand{\milcmarca}{\framebox[3.8mm]{\rule{0pt}{3.4mm}}}

% Fila marcable de lista suelta (checks de la estacion 1).
\newcommand{\milccheckrow}[1]{%
\par\vspace{1.8mm}\noindent
\makebox[6.4mm][l]{\raisebox{-.55mm}{\textcolor{milcNegro}{\milcmarca}}}%
\parbox[t]{\dimexpr\linewidth-6.4mm\relax}{\RaggedRight #1}\par}

% Celda marcable dentro de la rubrica (tabla de autoevaluacion). Misma
% sangria francesa, pero pensada para vivir dentro de una columna X.
\newcommand{\milcopcion}[1]{%
\makebox[5.2mm][l]{\raisebox{-.55mm}{\textcolor{milcNegro}{\milcmarca}}}%
\parbox[t]{\dimexpr\linewidth-5.2mm\relax}{\RaggedRight #1}}

% ── Recursos visuales de la guía ──────────────────────────────────────
% Declarados en el bloque `recursos:` del YAML y emitidos por el builder.
% \guiaFigura   → imagen rasterizada o PDF (foto, carta celeste, montaje).
% \guiaDiagrama → fuente TikZ/circuitikz (.tex) incrustada como vector:
%                 es la vía para circuitos y esquemáticos editables.
% [#1] = texto alternativo (alt) · #2 = ruta relativa al .tex · #3 = pie
% El alt llega al PDF etiquetado (/Alt) y lo lee el lector de pantalla.
\newcommand{\guiaFigura}[3][]{%
\begin{center}
\includegraphics[width=0.88\linewidth,height=0.38\textheight,keepaspectratio,alt={#1}]{#2}\\[1.5mm]
{\footnotesize\itshape\textcolor{milcNegro!75}{#3}}
\end{center}
}

\newcommand{\guiaDiagrama}[2]{%
\begin{center}
\input{#1}\\[1.5mm]
{\footnotesize\itshape\textcolor{milcNegro!75}{#2}}
\end{center}
}

\begin{document}
\thispagestyle{empty}

% ── Portada ───────────────────────────────────────────────────────────
% Antes: un rectangulo de color a pagina completa. Se veia bien en pantalla,
% pero en la fotocopiadora del colegio se comia el toner y salia gris sucio.
% Ahora la banda ocupa el tercio superior y el resto va en blanco.
\begin{tikzpicture}[remember picture,overlay]
  \path[fill=milcGradePrimary]
    (current page.north west) rectangle ([yshift=-10.6cm]current page.north east);

  \node[anchor=north west,text=milcGradeText] at ([xshift=2.0cm,yshift=-1.55cm]current page.north west)
    {\sffamily\bfseries\fontsize{12}{14}\selectfont GRADO};
  \node[anchor=north west,text=milcGradeText,opacity=.85] at ([xshift=2.0cm,yshift=-2.35cm]current page.north west)
    {\sffamily\bfseries\fontsize{8.5}{10}\selectfont SERIE GUÍAS · TECNOLOGÍA E INFORMÁTICA};
  \node[anchor=north east,text=milcGradeText,opacity=.85,align=right] at ([xshift=-2.0cm,yshift=-1.55cm]current page.north east)
    {\sffamily\bfseries\fontsize{9}{12}\selectfont I.E. Sor María Juliana\\Cartago, Valle del Cauca};

  \node[anchor=west,text=milcGradeText] (gradonum) at ([xshift=1.85cm,yshift=-7.1cm]current page.north west)
    {\sffamily\bfseries\fontsize{112}{112}\selectfont 9};
  % El circulito de grado (°) solo aplica a las guias de grado («11°»); en
  % Bebras y Semillero el numero grande es un momento/modulo, no un grado.
  % Se posiciona relativo al nodo `gradonum`, no en coordenadas absolutas.
  \draw[milcGradeText,line width=1.6pt]
    ([xshift=2.0mm,yshift=-4.5mm]gradonum.north east) circle (.15cm);

  \node[anchor=west,text=milcGradeText,text width=11.4cm,align=left] at ([xshift=1.5cm,yshift=0.5cm]gradonum.east)
    {
     {\sffamily\bfseries\fontsize{9}{11}\selectfont\MakeUppercase{Período 3 · Datos --- del registro al insight}}\\[3mm]
     {\sffamily\bfseries\fontsize{21}{25}\selectfont\setlength{\baselineskip}{27pt}Visualización honesta ---\\[4pt]gráficos\\[4pt]que mienten}\\[3.5mm]
     {\sffamily\bfseries\fontsize{15}{18}\selectfont Guía 28-9-TIC}
    };
\end{tikzpicture}

\vspace*{9.4cm}
\vfill

{\sffamily\bfseries\fontsize{8}{10}\selectfont\textcolor{milcGradeDark}{LO QUE TE LLEVAS HOY}}

\begin{tcolorbox}[enhanced,colback=white,colframe=milcNegro,arc=2mm,boxrule=1.6pt,
  left=5mm,right=5mm,top=4mm,bottom=4mm,before skip=3pt,after skip=12pt,
  fontupper=\RaggedRight\fontsize{11}{15.5}\selectfont]
La auditoría de tres gráficos reales tomados de prensa, redes o publicidad, cada uno con su captura y su fuente. De cada uno se nombra el truco usado de la lista de cinco, se explica con lo que se ve en la imagen y se dice qué decisión cambiaría si alguien lo creyera. Se cierra con la versión honesta de cada gráfico, hecha con los mismos datos.
\end{tcolorbox}

\begin{tcolorbox}[enhanced,colback=milcGris,colframe=milcGris,arc=2mm,boxrule=0pt,
  left=5mm,right=5mm,top=3.5mm,bottom=3.5mm,after skip=12pt]
{\sffamily\bfseries\fontsize{8}{10}\selectfont\textcolor{milcNegro!55}{ESTA GUÍA ES DE}}\\[3mm]
\begin{tabularx}{\linewidth}{X p{3.2cm} p{3.2cm}}
\linead{\linewidth} & \linead{3.0cm} & \linead{3.0cm}\\[-1mm]
{\fontsize{8.5}{10}\selectfont\textcolor{milcNegro!55}{Nombre completo}} &
{\fontsize{8.5}{10}\selectfont\textcolor{milcNegro!55}{Grupo}} &
{\fontsize{8.5}{10}\selectfont\textcolor{milcNegro!55}{Fecha}}\\
\end{tabularx}
\end{tcolorbox}

\vfill

\noindent\textcolor{milcLinea}{\rule{\linewidth}{1pt}}\\[2mm]
{\sffamily\bfseries\fontsize{9.5}{12}\selectfont Autor: Dr. Álvaro Cárdenas Orozco}\\[1mm]
{\sffamily\fontsize{8.5}{11}\selectfont\textcolor{milcNegro!55}{Metodología MILC · apertura ancestral · pensamiento computacional · triángulo Dussel–estoicismo–Floridi}}

\newpage

\section*{Apertura: Visualización honesta --- gráficos que mienten}

\begin{softbox}{milcGradeDark}{milcGradeSoft}{Saber ancestral}
Cada agosto Cartago celebra su aniversario, con misa en la catedral, ofrenda floral y sesión solemne del Concejo (Alcaldía de Cartago, 2021). Un detalle que casi nadie menciona: la fundación de 1540 ocurrió donde hoy está Pereira. A pocos kilómetros, Ansermanuevo celebra a mediados del mismo mes sus \textbf{Fiestas del Calado y el Bordado}. Se llama a sí misma cuna del calado y el bordado. Dos pueblos vecinos, dos ferias en el mismo mes, y en ambas la misma escena: en la feria uno se pone lo mejor que tiene. La pregunta incómoda viene después. ¿Lo mejor que tiene es lo que hace o lo que compra? No se trata de juzgar a nadie por cómo se presenta. Se trata de notar que presentarse es una decisión, y que la presentación puede coincidir con lo que hay detrás o puede taparlo. Un gráfico hace exactamente eso: presenta.\par\smallskip{\footnotesize\textcolor{milcNegro!70}{\textbf{Fuente.} Alcaldía de Cartago. (2021, 9 de agosto). \emph{Cartago celebra 481 años de fundación}. Municipio de Cartago Valle.}}
\end{softbox}

\begin{softbox}{milcTurquesa}{milcGris}{Saber contemporáneo}
Un gráfico engañoso casi nunca miente con los números. Miente con la selección y con la presentación, y por eso es difícil de atrapar. Hay cinco trucos que cubren casi todos los casos. \textbf{Eje cortado}: el eje vertical no arranca en cero y una diferencia pequeña parece enorme. \textbf{Escala torcida}: el espacio entre valores no es constante y aplana lo que crece. \textbf{Rango elegido}: se muestran solo los meses que convienen. \textbf{Categorías escondidas}: se dibujan tres barras y había siete. \textbf{Contexto omitido}: falta la cifra con la que habría que comparar. Y hay un procedimiento de diez segundos que atrapa la mayoría. Mira el eje vertical y comprueba si arranca en cero. Mira el rango de fechas. Cuenta las categorías visibles y pregúntate cuántas faltan. Y pregúntate qué cifra necesitarías para juzgar si eso es mucho o poco.
\end{softbox}

\begin{softbox}{milcMostaza}{milcGris}{Pregunta-puente}
\textit{En la feria uno se pone lo mejor que tiene, y eso no engaña a nadie porque todos saben que es una feria. ¿Cuándo un gráfico deja de presentarse y empieza a engañar?}
\end{softbox}

\begin{softbox}{milcVerde}{milcGris}{Saber y saber hacer}
Al terminar podrás: (1) \textbf{identificar} qué conclusión te empuja a sacar un gráfico antes de saber si es honesto; (2) \textbf{analizar} los cinco trucos y aplicar el chequeo de diez segundos; (3) \textbf{evaluar} tres gráficos reales, nombrar su truco y rehacerlos honestamente con los mismos datos.
\end{softbox}



\small
\begin{tabularx}{\textwidth}{>{\bfseries}p{3.0cm}Y}
\rowcolor{milcGradePrimary!12}Autor & Dr. Álvaro Cárdenas Orozco\\
DBA & Tratamiento y visualización honesta de datos para la toma de decisiones (MEN, T\&I)\\
Referentes & Ética del dato · Visualización honesta · Pensamiento computacional\\
Producto final & La auditoría de tres gráficos reales tomados de prensa, redes o publicidad, cada uno con su captura y su fuente. De cada uno se nombra el truco usado de la lista de cinco, se explica con lo que se ve en la imagen y se dice qué decisión cambiaría si alguien lo creyera. Se cierra con la versión honesta de cada gráfico, hecha con los mismos datos.\\
\end{tabularx}
\normalsize

\puente{En la sesión 6 hiciste gráficos y en la 7 limpiaste datos. Hoy miras el otro lado: qué pasa cuando alguien usa esas mismas herramientas para empujarte a una conclusión. En la sesión 9 conviertes datos en decisiones.}

\milcestacion{milcGradeDark}{Ruta MILC: así vamos a aprender}
\begin{tabularx}{\textwidth}{>{\centering\arraybackslash}X>{\centering\arraybackslash}X>{\centering\arraybackslash}X>{\centering\arraybackslash}X}
\phasecard{1}{milcVerde}{milcGris}{Escucha\\[-1mm]\small Observo, escucho y pregunto.} &
\phasecard{2}{milcTurquesa}{milcGris}{Entender\\[-1mm]\small Sistematizo y ordeno.} &
\phasecard{3}{milcMagenta}{milcGris}{Hacer\\[-1mm]\small Construyo y aplico.} &
\phasecard{4}{milcVino}{milcGris}{Cierre\\[-1mm]\small Evalúo y reflexiono.}
\end{tabularx}


\puente{Antes de la teoría vas a mirar tres gráficos sin que te digan que son engañosos. Lo que se te ocurra en esos primeros diez segundos es lo que le pasa a cualquier lector.}

\milcestacion{milcVerde}{Estación 1 de 3 · Escucha}

\begin{softbox}{milcVerde}{milcGris}{Escena para escuchar}
\textbf{Actividad 1 · IDENTIFICA --- Mira sin sospechar} (15 min · individual). (1) Tu docente proyecta tres gráficos reales, sin decir si tienen algo raro. (2) Para cada uno, escribe qué pregunta parece responder. (3) Escribe qué conclusión sacaste en los primeros diez segundos. (4) Ahora míralos otra vez con calma y anota si cambiarías esa conclusión. (5) Marca en cuál cambió más entre la primera y la segunda mirada.
\end{softbox}

{\sffamily\bfseries\fontsize{8}{10}\selectfont\textcolor{milcNegro!55}{MARCA CUANDO LO HAYAS HECHO}}
\milccheckrow{Anoté la conclusión de los primeros diez segundos en los tres.}
\milccheckrow{Volví a mirarlos con calma y anoté si cambió.}
\milccheckrow{Marqué en cuál cambió más.}

\begin{infoband}{milcVerde}


\textbf{Lo que va al cuaderno (Actividad 1 · IDENTIFICA).} \textbf{Título:} «Actividad 1 --- Mira sin sospechar». \textbf{Formato:} tabla de 3 filas y 3 columnas (gráfico / conclusión a los diez segundos / conclusión con calma), con marca en el que más cambió. \textbf{Extensión:} un tercio de página. \textbf{Sabes que terminaste cuando} las tres filas tienen las dos conclusiones escritas.
\end{infoband}

\puente{Ya viste tu propia reacción. Ahora vas a nombrar los cinco trucos y a aprender el chequeo que los atrapa en diez segundos.}

\milcestacion{milcTurquesa}{Estación 2 de 3 · Entender}

\begin{softbox}{milcTurquesa}{milcGris}{Lo que vas a entender}
\textbf{Actividad 2 · ANALIZA --- Los cinco trucos} (20 min · parejas). (1) Con tu pareja, escriban los cinco trucos con una frase propia cada uno. (2) Vuelvan a los tres gráficos y busquen cuál truco usa cada uno. (3) Apliquen el chequeo de diez segundos y anoten qué encontró cada paso. (4) Escriban qué decisión cambiaría si alguien creyera el gráfico tal como está.
\end{softbox}

\milcpaso{1}{milcTurquesa}{\textbf{Los cinco trucos.} Eje cortado, escala torcida, rango elegido, categorías escondidas y contexto omitido. Casi todos los gráficos engañosos usan uno de esos cinco, y muchos usan dos a la vez.}
\milcpaso{2}{milcTurquesa}{\textbf{El chequeo de diez segundos.} Mira si el eje vertical arranca en cero. Mira el rango de fechas. Cuenta las categorías visibles. Pregúntate qué cifra falta para saber si eso es mucho o poco.}
\milcpaso{3}{milcTurquesa}{\textbf{Casi nunca mienten con los números.} Mienten con la selección y con la presentación. Por eso no sirve pedir que «pongan datos reales»: los datos suelen ser reales y aun así el gráfico empuja a una conclusión falsa.}
\milcpaso{4}{milcTurquesa}{\textbf{Se demuestra señalando, no sospechando.} «Este gráfico se ve raro» no es una auditoría. «El eje arranca en 48 y no en cero, y por eso una diferencia de dos puntos ocupa media pantalla» sí lo es.}

\begin{drawbox}{milcTurquesa}{El chequeo de diez segundos}
\textbf{1. Eje vertical:} ¿arranca en cero? \\[1mm] \textbf{2. Fechas:} ¿qué rango muestra y cuál omite? \\[1mm] \textbf{3. Categorías:} ¿cuántas se ven y cuántas había? \\[1mm] \textbf{4. Contexto:} ¿con qué cifra habría que comparar?
\end{drawbox}

\begin{softbox}{milcMostaza}{milcGris}{Errores comunes y su corrección}
(1) \textbf{Auditar por sospecha}: sin señalar qué elemento produce el efecto, no hay argumento. (2) \textbf{Creer que el truco está en los números}: casi siempre está en la selección. (3) \textbf{Olvidar el rango de fechas}: es el truco más común y el que menos se mira. (4) \textbf{Confundir feo con engañoso}: un gráfico mal hecho puede ser honesto. (5) \textbf{Rehacerlo con otros datos}: la versión honesta usa los mismos datos, o no demuestra nada.
\end{softbox}

\begin{infoband}{milcTurquesa}


\textbf{Lo que va al cuaderno (Actividad 2 · ANALIZA).} \textbf{Título:} «Actividad 2 --- Los cinco trucos». \textbf{Formato:} los cinco trucos con frase propia, el truco de cada uno de los tres gráficos y los cuatro pasos del chequeo con lo que encontró cada uno. \textbf{Extensión:} media página. \textbf{Sabes que terminaste cuando} cada truco identificado señala un elemento visible del gráfico.
\end{infoband}

\puente{Tienes el método. Toca salir a buscar gráficos reales y rehacerlos honestamente con sus propios datos.}

\milcestacion{milcMagenta}{Estación 3 de 3 · Hacer}

\begin{softbox}{milcMagenta}{milcGris}{Cómo vas a construir tu producto}
\textbf{Actividad 3 · EVALÚA --- Auditoría de tres gráficos} (30 min · individual). (1) Busca tres gráficos reales en prensa, redes, publicidad o un libro. (2) Guarda la captura y la fuente de cada uno. (3) Nombra el truco de cada uno usando la lista de cinco. (4) Explica el truco señalando qué se ve en la imagen, con números. (5) Escribe qué decisión cambiaría si alguien lo creyera. (6) Rehaz cada gráfico de forma honesta, con los mismos datos, y ponlos lado a lado.
\end{softbox}

\milcpaso{1}{milcMagenta}{\textbf{Tu entrega tiene tres piezas por gráfico.} La captura con su fuente. La auditoría con el truco señalado. Y la versión honesta hecha con los mismos datos.}
\milcpaso{2}{milcMagenta}{\textbf{La fuente es obligatoria.} Un gráfico sin origen no se puede auditar ni discutir. Si no encuentras de dónde salió, ya tienes un hallazgo, pero busca otro para auditar.}
\milcpaso{3}{milcMagenta}{\textbf{La versión honesta usa los mismos datos.} Si cambias las cifras, no demostraste nada. Lo que cambias es el eje, el rango, las categorías visibles o el contexto que añades.}
\milcpaso{4}{milcMagenta}{\textbf{Di qué decisión cambiaría.} Es lo que separa un ejercicio de una auditoría. «Alguien podría creer que el problema se duplicó cuando subió dos puntos» tiene consecuencias; «está mal hecho» no las tiene.}

\begin{drawbox}{milcMagenta}{Antes de entregar}
\checkbox Tres gráficos reales con captura y fuente verificable. \\[1mm] \checkbox Tres trucos nombrados con la lista de cinco. \\[1mm] \checkbox Cada truco señalado en la imagen, con números. \\[1mm] \checkbox Tres versiones honestas con los mismos datos. \\[1mm] \checkbox Cada uno dice qué decisión cambiaría. \\[1mm] \checkbox Ningún argumento se apoya solo en que «se ve raro».
\end{drawbox}

\begin{softbox}{milcMagenta}{milcGris}{Plantilla de trabajo}
\textbf{Gráfico.} \textit{Fuente y fecha.} \\[1mm] \textbf{Truco.} \textit{Cuál de los cinco.} \\[1mm] \textbf{Prueba.} \textit{Qué se ve en la imagen, con números.} \\[1mm] \textbf{Consecuencia.} \textit{Qué decisión cambiaría.} \\[1mm] \textbf{Versión honesta.} \textit{Mismos datos, qué cambié.}
\end{softbox}

\begin{infoband}{milcMagenta}


\textbf{Lo que va al cuaderno (Actividad 3 · EVALÚA).} \textbf{Título:} «Actividad 3 --- Auditoría de tres gráficos». \textbf{Formato:} los tres con fuente, truco, prueba señalada, consecuencia y el boceto de la versión honesta. \textbf{Extensión:} una página. \textbf{Sabes que terminaste cuando} cada truco está probado con un número de la imagen.
\end{infoband}

\puente{Lo que entregas hoy: tres auditorías y tres versiones honestas. La prueba de calidad: cada truco se demuestra señalando la imagen, no sospechando.}

\milcestacion{milcMostaza}{Producto de la sesión}
\textbf{Tres gráficos auditados con su truco probado + tres versiones honestas}.
\begin{softbox}{milcMostaza}{milcGris}{Criterios de calidad}
(1) Los tres gráficos son reales y tienen captura y fuente verificable. (2) Cada truco está nombrado con la lista de cinco. (3) Cada truco se prueba señalando qué se ve en la imagen, con números. (4) Las tres versiones honestas usan los mismos datos del original. (5) Cada auditoría dice qué decisión cambiaría si alguien creyera el gráfico.
\end{softbox}

\puente{Antes de cerrar, mira lo que hiciste desde cinco lados. Un gráfico bien hecho respeta al lector; uno torcido cuenta con que no va a mirar el eje.}

\milcestacion{milcVino}{Evaluación liberadora · cinco dimensiones}

\begin{softbox}{milcVino}{milcGris}{Sentido de la evaluación}
Evaluar no es solo revisar una nota. Es reconocer qué aprendí, qué cambió en mí, qué emociones manejé, cómo me relaciono con la ciudad y la comunidad, y qué le dejaré a quien viene detrás.
\end{softbox}

\textbf{Mi reflexión liberadora en cinco dimensiones.} Completa cada frase iniciada:

\small
\begin{tabularx}{\textwidth}{>{\bfseries\RaggedRight\arraybackslash}p{3.4cm}Y>{\RaggedRight\arraybackslash}p{4.6cm}}
\rowcolor{milcVino}\color{white}Dimensión & \color{white}\bfseries Mi respuesta (frame starter) & \color{white}\bfseries Algo que descubrí\\
1. Personal & Lo que más cambió en mí fue \linead{6.5cm} & \linead{4.2cm}\\
2. Emocional & La emoción más fuerte fue \linead{2.5cm} y la regulé \linead{3cm} & \linead{4.2cm}\\
3. Ciudadana & En democracia esto importa porque \linead{6cm} & \linead{4.2cm}\\
4. Local (Cartago) & En mi barrio sería útil para \linead{6.5cm} & \linead{4.2cm}\\
5. Intergeneracional & A mi abuela/abuelo le diría \linead{6.5cm} & \linead{4.2cm}\\
\end{tabularx}
\normalsize

\begin{infoband}{milcVino}
\textbf{Para la próxima sustentación.} Vuelve a esta tabla en seis meses. Verás que las respuestas cambian — no porque lo aprendido cambie, sino porque tú cambias. La evaluación liberadora es una conversación contigo a través del tiempo.
\end{infoband}


\puente{Para terminar, tres citas verificadas sobre presentar y aparentar: Dussel sobre una estética más allá de la publicidad, Marco Aurelio sobre ser en vez de discutirlo, Floridi sobre la abundancia que distrae.}

\milcestacion{milcGradeDark}{Triángulo de pensamiento · cierre filosófico}

\begin{softbox}{milcVino}{milcGris}{Dussel · lente del nosotros}
\textit{«El arte popular es el arte primero, la suprema expresión de la estética. Se da en la vida cotidiana, en la música, la danza, la pintura, en el teatro… Es necesario formular una estética popular más allá de la publicidad y la moda.»}

{\footnotesize\itshape\color{milcNegro!70}--- Enrique Dussel · Filosofía de la liberación (1977), §4.2.9.3}

Un gráfico también es una pieza estética, y por eso convence antes de que uno lo lea. La diferencia entre presentar y persuadir es la misma que hay entre ponerse lo mejor que uno tiene y aparentar lo que no se tiene.

\textbf{¿Cuándo he elegido un gráfico porque se veía impresionante y no porque mostrara mejor el dato?}
\end{softbox}
\linead{\linewidth}\\[1mm]
\linead{\linewidth}

\begin{softbox}{milcMostaza}{milcGris}{Estoicismo · Marco Aurelio · Meditaciones X, 16 (c. 175 d.C.) · cuidado interior}
\textit{«De hoy más, déjate absolutamente de disputar cuál conviene que sea un hombre bueno, sino procura ser tal en realidad.»}

{\footnotesize\itshape\color{milcNegro!70}--- Marco Aurelio · Meditaciones X, 16 (c. 175 d.C.)}

Con los gráficos igual. No se trata de escribir que el tuyo es honesto: se trata de que el eje arranque en cero y de que el rango de fechas sea el completo. La honestidad de un gráfico se ve en el gráfico.

\textbf{¿Qué elemento de mi último gráfico no resistiría el chequeo de diez segundos?}
\end{softbox}
\linead{\linewidth}\\[1mm]
\linead{\linewidth}

\begin{softbox}{milcTurquesa}{milcGris}{Floridi · lente de la infoesfera}
\textit{«La abundancia de información también puede producir sobrecarga cognitiva, distracción y amnesia (el presente olvidadizo). (trad. propia)»}

{\footnotesize\itshape\color{milcNegro!70}--- The Onlife Initiative (ed. Luciano Floridi) · The Onlife Manifesto (2015), § 2.3}

Los gráficos engañosos viven de eso. Nadie tiene tiempo de auditar cada imagen que pasa, y el truco cuenta con esa prisa. Por eso el chequeo tiene que caber en diez segundos: si no cabe, nadie lo va a hacer.

\textbf{¿Cuántos gráficos vi esta semana, y a cuántos les miré el eje?}
\end{softbox}
\linead{\linewidth}\\[1mm]
\linead{\linewidth}

\begin{infoband}{milcNegro}
\textbf{Nota para el docente.} Las tres son citas textuales y llevan su referencia debajo. Puedes remitir a la obra en clase.
\end{infoband}

\milcestacion{milcVino}{Mi autoevaluación · marca una casilla por fila}

{\small
\setlength{\extrarowheight}{2.2pt}
\begin{tabularx}{\textwidth}{>{\RaggedRight\arraybackslash\bfseries}p{3.0cm}YYY}
\rowcolor{milcVino}\color{white}\bfseries Criterio & \color{white}\bfseries Lo logré & \color{white}\bfseries En proceso & \color{white}\bfseries Necesito apoyo\\[.6mm]
Comprensión del tema & \milcopcion{Explico las ideas con mis palabras.} & \milcopcion{Reconozco las ideas, falta precisión.} & \milcopcion{Necesito repasar lo básico.}\\[.8mm]
\rowcolor{milcGris}Pensamiento computacional & \milcopcion{Usé los pilares al armar mi producto.} & \milcopcion{Usé algunos pilares.} & \milcopcion{Necesito releer la estación 2.}\\[.8mm]
Aplicación y producto & \milcopcion{Mi producto cumple los criterios.} & \milcopcion{Está armado, con ajustes pendientes.} & \milcopcion{Aún está en construcción.}\\[.8mm]
\rowcolor{milcGris}Reflexión 5 dimensiones & \milcopcion{Respondí cada una con honestidad.} & \milcopcion{Respondí algunas con detalle.} & \milcopcion{Mis respuestas son muy generales.}\\[.8mm]
Triángulo de pensamiento & \milcopcion{Conecté las 3 voces con mi trabajo.} & \milcopcion{Conecté al menos una.} & \milcopcion{No las relaciono con mi proceso.}\\[.8mm]
\end{tabularx}}

\vspace{3mm}
\textbf{Hoy entendí que:}\\[1mm]
\linead{\linewidth}\\[1mm]
\linead{\linewidth}

\textbf{Mi compromiso para la próxima sesión es:}\\[1mm]
\linead{\linewidth}\\[1mm]
\linead{\linewidth}

\begin{softbox}{milcMostaza}{milcGris}{Cita de despedida · Marco Aurelio}
\textit{``El obstáculo en el camino se vuelve el camino. Lo que estorba al camino se vuelve camino.''}\\[1mm]
Cada error de hoy, bien observado, se convierte en pieza del camino que pisarás mañana.
\end{softbox}

\small
\begin{tabularx}{\textwidth}{>{\bfseries}p{3.6cm}Y}
\rowcolor{milcGradeDark!12}Nombre completo: & \linead{10cm}\\
Fecha: & \linead{4cm} \quad Curso/grupo: \linead{4cm}\\
Mi firma: & \linead{9cm}\\
\end{tabularx}
\normalsize

\begin{infoband}{milcGradeDark}
\textbf{Para la próxima sesión.} Llega con las tres auditorías y las versiones honestas. Cuando veas un gráfico esta semana, hazle el chequeo de diez segundos antes de creerlo. En la sesión 9 conviertes tus propios datos en decisiones.
\end{infoband}

\end{document}
